\section{Graded rings and the Proj construction}

\begin{definition} \label{def:graded ring}
In what follows, let $P := \bigoplus_{d \geq 0} P_d$ be a fixed graded ring satisfying the following conditions.

\begin{itemize}
\item[(i)]
The subring $P_0$ of $P$ is a field and $P_1 \neq 0$.
\item[(ii)]
The multiplicative monoid 
\[
\left( \bigcup_{n=0}^\infty P_n \setminus \{0\} \right)/P_0^\times
\]
is generated by $(P_1 \setminus \{0\})/P_{0}^{\times}$.
\item[(iii)]
For each nonzero $s \in P_1$, there exists a graded ring isomorphism 
\begin{equation} \label{eq:graded ring isom mod s}
P/sP \cong P_0 \oplus T L_s[T]
\end{equation}
for some field $L_s$ containing $P_0$. 
\end{itemize}

\end{definition}

\begin{lemma} \label{prop:PID from graded ring}
\uses{def:graded ring}
For $P$ as in Definition~\ref{def:graded ring}, for each nonzero $s \in P_1$, $P[s^{-1}]_0$
is a principal ideal domain. 
\end{lemma}
\begin{proof}
\end{proof}

\begin{definition} \label{def:line bundles}
\uses{def:graded ring,def:vector bundles}
For $P$ as in Definition~\ref{def:graded ring},
for $d \in \Z$, define the line bundle $\cO_X(d)$ on $X$ as $\Proj \bigoplus_{n=0}^\infty P_{n+d}$.
\end{definition}

\begin{lemma} \label{lem:dimension at least 2}
\uses{def:graded ring}
For $P$ as in Definition~\ref{def:graded ring}, we have $\dim_{P_0} P_1 \geq 2$.
In particular, there exist $s,t \in P_1$ which are linearly independent over $P_0$.
\end{lemma}
\begin{proof}
By condition (i) of Definition~\ref{def:graded ring},
we can pick a nonzero $s \in P_1$.
By condition (iii) of Definition~\ref{def:graded ring}, we can find $t \in P_1$ 
with nonzero image in $P_1/sP_0$. This proves the claim.
\end{proof}

\begin{lemma} \label{lem:H1 of line bundles}
\uses{def:line bundles}
We have $H^1(X, \cO_X) = 0$.
\end{lemma}
\begin{proof}
\uses{lem:dimension at least 2}
For any $s,t \in P_1$ which are linearly independent over $P_0$
(as in Lemma~\ref{lem:dimension at least 2},
by \eqref{eq:graded ring isom mod s} we have $P_2 = P_1 s + P_1 t$.
By induction on $m+n+k$ we deduce that $P_{m+n+k} = P_{m+k} s^n + P_{n+k} t^m$ for any $m,n >0$ and $k \geq 0$, and hence
\[
H^1(X, \cO_X) = \frac{P[(st)^{-1}]_0}{P[s^{-1}]_0 + P[t^{-1}]_0} = 
\varinjlim_m \frac{P_{2m}}{P_{m}s^m + P_{m} t^m} = 0. 
\qedhere
\]
\end{proof}

\begin{proposition} \label{prop:graded ring to curve}
\uses{def:graded ring,def:abstract complete curve}
For $P$ as in Definition~\ref{def:graded ring}, $X := \Proj(P)$ is an abstract complete curve
for the degree function carrying every closed point to $1$.
\end{proposition}
\begin{proof}
\uses{prop:PID from graded ring,lem:dimension at least 2}
Pick $s,t$ as in Lemma~\ref{lem:dimension at least 2};
by Lemma~\ref{lem:H1 of line bundles},
$X$ is covered by the distinguished open subspaces corresponding to $s$ and $t$.
By Proposition~\ref{prop:PID from graded ring},  each of these subspaces is regular of dimension 1.
By condition (iii) of Definition~\ref{def:graded ring}, each closed point of either subspace has algebraically closed residue field.
\end{proof}

\begin{proposition} \label{prop:Pic of curve}
\uses{prop:graded ring to curve,def:line bundles}
The degree map $\deg\colon \Div(X) \to \Z$ induces an isomorphism $\Pic(X) \cong \Z$.
\end{proposition}
\begin{proof}
\uses{prop:PID from graded ring}
The map $\Pic(X) \to \Z$ carries $\cO_X(n)$ to $n$, so we need only check injectivity.
This follows from Proposition~\ref{prop:PID from graded ring}.
\end{proof}

\begin{theorem} \label{thm:graded ring projective line}
\uses{def:graded ring,def:generalized projective line}
For $P$ as in Definition~\ref{def:graded ring}, $X := \Proj(P)$ is a generalized projective line.
In particular the conclusion of Theorem~\ref{thm:abstract classification} applies.
\end{theorem}
\begin{proof}
\uses{prop:PID from graded ring,prop:Pic of curve,lem:dimension at least 2,lem:H1 of line bundles}
By Proposition~\ref{prop:graded ring to curve}, $X$ is an abstract complete curve.
Of the conditions from Definition~\ref{def:generalized projective line}, 
(i) comes from Proposition~\ref{prop:Pic of curve};
(ii) comes from Lemma~\ref{lem:H1 of line bundles}; and
(iii) comes from the description of the degree function in Proposition~\ref{prop:graded ring to curve}.
\end{proof}